\documentclass[10pt]{article}
\usepackage[margin=20mm]{geometry}
\usepackage{amsmath,amssymb,mathtools,fontspec,unicode-math,xurl,hyperref}
\IfFontExistsTF{DejaVu Serif}{\setmainfont{DejaVu Serif}}{\setmainfont{Latin Modern Roman}}
\setmathfont{latinmodern-math.otf}
\setlength{\parindent}{0pt}\setlength{\parskip}{6pt}\sloppy
\title{Technical verification of New Garside structures and applications to Artin groups}
\author{AI-assisted technical review}
\date{}
\begin{document}
\maketitle
Abstract: Three statements are false with their printed hypotheses: the general interval-label criterion, the restricted systolic-presentation theorem, and the positive-square theorem. Explicit torsion examples establish these failures. The special positive-germ construction has stronger properties than the defective background definition; the Artin-group conclusions are not disproved by these examples. Local corrections to the lattice and support arguments, a checked unmarked complex identification in the homotopy argument, exact finite Coxeter checks, and the remaining standard-theorem dependencies are given below. Scientific adjudication remains pending.\par
\section*{Exact source, scope and assessment}
The reviewed object is Thomas Haettel and Jingyin Huang, New Garside structures and applications to Artin groups, M25-0168, arXiv:2305.11622v2, stamped 16 October 2024 and internally dated 17 October 2024, 50 pages [source]. The assignment associates it with DOI 10.1215/00127094-2024-0062 (https://doi.org/10.1215/00127094-2024-0062). Identity with the journal-final article or a subsequent correction is not asserted. Page numbers below are one-based physical PDF pages, agreeing with printed numbers.\par
The lead read the complete supplied text, including all definitions, proofs, examples and references. Three native helpers independently checked Section 3, Section 5, and Section 6 with their dependencies. The lead actually viewed 31 source pages; combined actual source-image coverage is all 50 pages. The checks distinguish simple-generator length, reflection length, partial-product divisibility and ambient absolute order, and include identity cases, empty supports, repeated letters, parabolic intersections and the scope of geometric consequences.\par
The supported impact is Major: Theorems 3.13 and 3.15 admit explicit counterexamples under their printed hypotheses. The overbroad background criterion, Theorem 2.9, also has a counterexample, but its effect is Minor here: the actual positive-germ construction has the additional nondegeneracy and interval structure needed for the valid standard theorem. This does not assert that the Artin-group main theorems are false. Their positive-germ route survives these examples, subject to the repairs and scoped standard dependencies below.\par
\section*{The arbitrary interval-label criterion}
Problem 1 (Minor background qualification; Definitions 2.6–2.8 and Theorem 2.9, pp.15–17).  A group-like balanced labeling of a finite graded lattice, in the literal sense defined here, need not have a Garside interval group.\par
Take \(P=\{0<1<2<3\}\), with \[\lambda(i,j)=j-i\pmod 2\quad\hbox{in }C_2=\{z,x\},\] where \(z\) is the identity label. This is a bounded finite graded lattice. Labels on every weak three-point chain multiply in \(C_2\); cancellation makes any two corresponding labels determine the third. The bottom-label, top-label and complete label sets all equal \(\{z,x\}\). Both strengthened balanced equalities therefore hold.\par
Definition 2.7 identifies all vertices. The edge labels are \(z,x\), and the four strict three-point chains give \[x^2=z,\qquad xz=x,\qquad zx=x.\] The top three-cell does not change the fundamental group. Hence \[G_P=\langle x,z\mid x^2=z,\ xz=x,\ zx=x\rangle=C_2.\] Separately collapsing the identity-labeled edge leaves the same group. A nontrivial finite group cannot satisfy Definition 2.5: a nonidentity generator has finite order and yields arbitrarily long positive expressions for the identity, contrary to condition (2).\par
The original framework extends edge labels to languages of edge words in intervals [mcc05]. The arbitrary map in this manuscript discards that information. The assertion that every labeling in Definition 2.8 defines the required Noetherian germ is consequently false. Moreover, Digne’s Proposition 2.4(ii) requires cancellation against arbitrary words of the universal monoid, not just one local product [digne]. A correction must restore a nondegenerate interval-language/germ framework and its global theorem. Repeating the balanced equalities cannot suffice.\par
\section*{The systolic example theorem}
Problem 2 (Major; Definition 3.12, Theorem 3.13 and Corollary C, pp.2–3 and 28–29).  The two printed restricted-systolic conditions allow torsion.\par
Consider \[G=\langle a,b,c\mid aa=b,\ bb=c,\ cc=a\rangle.\] Eliminating \(b=a^2,c=a^4\) gives \(G=C_7\). The images of \(S=\{a,b,c\}\) are \(\{1,2,4\}\) modulo seven; their positives and negatives are all six nonzero residues, without repetitions. The unoriented Cayley graph is \(K_7\), without loops or multiple edges. Its flag completion is a six-simplex, hence simply connected, and every vertex link is a simplex and is 6-large. All defining relators have the required form \(st r^{-1}\), with repeated letters permitted. Both printed hypotheses hold.\par
But \(G\times\mathbb Z\) has order-seven torsion and cannot be Garside. For the standard torsion-free consequence, left divisibility extends to a lattice on a Garside fraction group. If \(g^m=1\), left multiplication by \(g\) permutes \(\{1,g,\ldots,g^{m-1}\}\); it fixes their join \(j\), and \(gj=j\) implies \(g=1\).\par
The construction fails explicitly at partial associativity: in \(U=\{1,a,b,c\}\), \((a*a)*b=b*b=c\) exists, whereas \(a*b=a^3\) is absent. Associativity of an ambient group does not imply Definition 3.1’s strong existence requirements. The genuine restricted-triangular framework includes \(stu\in S\Rightarrow st,tu\in S\) and its other presentation and local conditions [soergel]. Restoring those hypotheses excludes the example and narrows the statement. Torsion-freeness alone does not fix the proof: the free presentation \(\langle a,b,c,p,q\mid ab=p,\ pc=q\rangle\) has a systolic two-triangle flag complex, but \((a*b)*c\) exists while \(b*c\) does not. This latter observation refutes the germ inference, not the Garside-product conclusion for a free group.\par
\section*{The positive-square theorem}
Problem 3 (Major; Definition 3.14 and Theorem 3.15, pp.29–30).  The three link exclusions allow positive-square groups with torsion.\par
Use distinct \(a,b,c,d,e,f,g,h,x,y\) and eight further letters \(t_v\), \(v\in V=\{a,b,c,d,e,f,g,h\}\). Put \(H(v)=x\) on \(\{a,d,e,f\}\), and \(H(v)=y\) on \(\{b,c,g,h\}\), and impose \[ab=cd,\qquad ef=gh,\qquad t_vv=t_vH(v)\quad(v\in V).\] Repeated letters are expressly allowed in Definition 3.14. Cancellation and Tietze elimination give \[G\cong\langle x,y\mid xy=yx,\ x^2=y^2\rangle*F_8
   \cong(\mathbb Z\oplus C_2)*F_8.\] The element \(xy^{-1}\) has order exactly two: its nontriviality is detected by \(x\mapsto1,y\mapsto0\) in \(C_2\).\par
For \(pq=rs\), the link corner edges, with multiplicity retained, are \[\{p^+,r^+\},\quad\{q^-,s^-\},\quad
 \{p^-,q^+\},\quad\{r^-,s^+\}.\] The two main relations give outgoing edges \(a^+c^+,e^+g^+\), incoming edges \(b^-d^-,f^-h^-\), and mixed edges \(a^-b^+,c^-d^+,e^-f^+,g^-h^+\). A clone gives a loop at \(t_v^+\), an incoming edge \(v^-H(v)^-\), and mixed edges \(t_v^-v^+,t_v^-H(v)^+\). There are no parallel nonloop outgoing or incoming pairs. The endpoints of either outgoing nonloop edge have no common incoming neighbor. For every incoming edge, the endpoints have no common outgoing neighbor: the hubs \(x^-,y^-\) and the endpoints of the two main incoming edges have no mixed neighbors. The mixed graph is a forest of subdivided stars with incoming leaves attached. It contains no alternating embedded four-cycle. The loops are not any expressly excluded type. Exact independent enumeration retains all 40 corners and eight loops and finds zero forbidden patterns. Thus all hypotheses hold, whereas \(G\times\mathbb Z\) has torsion.\par
Banning loops or non-reduced relators alone is insufficient. For each clone introduce five fresh \(p_v,q_v,r_v,s_v,t_v\) and replace its equation by \[p_vv=q_vr_v,\qquad q_vr_v=s_vt_v,\qquad s_vt_v=p_vH(v).\] These equations give \(v=H(v)\); eliminating \(r_v=q_v^{-1}p_vv\) and \(t_v=s_v^{-1}p_vv\) leaves \(p_v,q_v,s_v\) free. The group becomes \((\mathbb Z\oplus C_2)*F_{24}\). Each gadget has an outgoing triangle on \(p_v^+,q_v^+,s_v^+\), an incoming path \(v^-,r_v^-,t_v^-,H(v)^-\), and mixed edges \(p_v^-v^+,p_v^-H(v)^+\) together with duplicated \(q_v^-r_v^+\) and \(s_v^-t_v^+\). Only mixed pairs are parallel, a type the theorem permits. The mixed simple graph is still a forest; same-type edges have no common opposite-type neighbor. All relators are cyclically reduced and no link loop occurs. An independent enumeration of 50 generators, 26 relations and 104 corners again gives zero forbidden patterns.\par
A sufficient narrowed formulation would require every nontrivial length-two equivalence class to be an isolated pair \(\{ab,cd\}\), with different first and last letters and no side shared by another relation, together with the printed link conditions. The resulting three-level germ has no nonidentity triple products; cancellation, atomic joins and the rectangle condition follow directly from the excluded configurations. This is a sufficient restriction, not a proposed characterization. The examples expose a substantive equivalence-chain obstruction.\par
\section*{The general construction and valid examples}
The doubled poset in Section 3 has properties absent from the counterexample in Section 2. Write its elements \(u_0,\bar u_1\). Its minimum is \(e_0\), its maximum is \(\Delta=\bar e_1\), and \(u_0\le\bar v_1\) means \(v*u\) exists. Both partial-associativity implications prove transitivity; positivity and cancellation prove antisymmetry. The identity case \(U=\{e\}\) is allowed. The three possible label relations give inverse maps \[G_E\longleftrightarrow G_U\times\mathbb Z,\qquad
 u_0\mapsto(u,0),\quad\bar u_1\mapsto(u^{-1},1).\] Indeed \(u_0\bar u_1=\Delta=\bar u_1u_0\), so \(\Delta\) is central and all upper generators equal \(u_0^{-1}\Delta\). Proposition 3.6 does not require a prior embedding of \(U\).\par
Proposition 3.7 incorrectly treats an arbitrary upper bowtie endpoint as minimal. A complete local repair is available. For lower \(u_0,v_0\) and upper \(w_0,\bar a_1\), take \(j=u\vee_Lv\le_Lw\); assumption (3) gives \(a*j\), so \(j_0\) is intermediate. The reversed distribution uses a right join and assumption (4), with the order reversed in the upper copy. For two endpoints in each copy, the rectangle hypothesis gives one of the needed joins and the multiplication hypotheses make it intermediate. Pure-copy cases use the corresponding meet. These exhaust the distributions. Bounded height ensures that a pair with an upper bound in a meet-semilattice has a join: a minimal upper bound \(j\) has \(j\wedge z=j\) for every other upper bound \(z\).\par
For this specific \(E\), \(\lambda(e_0,z)=z\), only degenerate intervals have identity label, and translation gives a label-preserving order isomorphism \[[x,\Delta]\longrightarrow[e_0,\partial x],\qquad
 y\mapsto\lambda(x,y),\qquad
 \partial u_0=\bar u_1,\quad\partial\bar u_1=u_0.\] Both directions follow from strong associativity and cancellation. All interval factorizations and atom-word languages are transported. Writing \(x\backslash y=\lambda(x,x\vee y)\) and \(J=x\vee y\vee z\), join transport gives \[(x\backslash y)\backslash(x\backslash z)
 =\lambda(x\vee y,J)
 =(y\backslash x)\backslash(y\backslash z).\] The opposite multiplication has the mirror identity. Positivity, nondegeneracy, bounded factorization height and these translations place \(E\) within the genuine interval-germ framework. Global cancellation, normal forms, fractions and asphericity remain the standard theorem dependency [mcc05, digne]; local cancellation is not a substitute. With this qualification, no counterexample to Theorems 3.8–3.9 is established.\par
In the \(T(5)\) branch, proper subwords are not newly identified and full-side classes cannot extend, proving strong associativity. Different first and last letters give cancellation. Bowties, multiplication of joins and rectangles respectively produce forbidden parallel corners, triangles and alternating squares. Identity, equality and comparability cases must be separated before asserting distinct signed vertices. Unused generators are adjoined as isolated atoms, including the zero-relator free-group case. The source’s sentence about proper subwords must exempt literally identical words. For the surface examples, elimination recovers the standard relations; their links have lengths \(6g-4\) for orientable \(g\ge2\), and \(2g\) for nonorientable \(g\ge3\). The sphere, torus, Klein bottle and free boundary cases also satisfy the conclusion; the excluded projective plane has torsion. These checks support Theorem 3.10 and Corollary 3.11.\par
\section*{Reflection germs and the Salvetti comparison}
The sentence after Lemma 4.1 identifying induced orders with restricted absolute order needs additional closure. In \(W(A_2)=\langle s,t\rangle\), let \(U=\{1,s,t,st\}\). The displayed factor conditions using \(R_U=\{s,t\}\) hold, but \(t\le_{\rm abs}st\) has no partial complement in \(U\), since \(tst\notin U\). All actual Artin applications use unions of full compatible spherical intervals and possess the needed closure. The general statement should use germ orders throughout, or explicitly impose that closure.\par
In Proposition 4.2, joins of covers \(vr_1,vr_2\) arise from \(r_1\vee r_2\) and atomic multiplication compatibility. The nested induction extending multiplication to longer elements works because a nonidentity remainder strictly increases the outer parameter. For the rectangle induction, \(a=e\) or \(b=e\) already gives a right join, and \(u=e\) or \(v=e\) a left join. Otherwise the positive remainder gives \(|a'x|>|x|\). The concluding \(aw=wu_1w''\) should be \(aw=au_1w''\). At the start of Section 4, reflection decomposition means a product in \(R\), not \(S\).\par
The proposed objection to Theorem 4.3’s complex identification is withdrawn. The natural marked poset map \(\bar u_1\mapsto u^{-1}\delta_Ts_0\) twists upper labels by conjugation, but page 34 asserts coincidence of the complexes, without prescribing this marking or the image of the top loop. There is an explicit unmarked cellular isomorphism.\par
Write \(U=U_T\), and let \(z\) be the positive generator of the new infinite cyclic Artin factor. The product dual interval is \(D=U\times\{1,z\}\); \(z\) projects to the Coxeter involution \(s_0\), but is not itself an involution in the interval group. On one-cells send \[u_0\longmapsto u,\qquad
 \bar v_1\longmapsto(vz)^{-1}.\] The second map reverses the orientation of the edge labeled \(vz\). In particular \(\Delta_E=\bar e_1\) maps to \(z^{-1}\), not to \(\delta_Tz\); the alleged centrality obstruction disappears.\par
For a mixed chain \[u_{0,0}<\cdots<u_{r,0}<\bar v_{0,1}<\cdots<\bar v_{s,1},\] the group vertices are \(u_i\) and \(v_j^{-1}z^{-1}\). Left translate by \(v_0z\) and cyclically reorder the vertices. The result is the chain in \(D\) \[1=v_0v_0^{-1}<\cdots<v_0v_s^{-1}
       <(v_0u_0)z<\cdots<(v_0u_r)z .\] The lower terms and upper terms belong to \(U\) because \(v_j\le_Rv_0\) and \(v_0*u_i\) exists. Their cross differences are \(v_j*u_i\), which are defined by the mixed-chain conditions. Consequently this is a simplex of the product dual interval.\par
Conversely a mixed chain \(a_0<\cdots<a_r<b_0z<\cdots<b_sz\) in \(D\) gives the chain in \(E_T\) \[(b_0^{-1}b_0)_0<\cdots<(b_0^{-1}b_s)_0
 <\overline{a_0^{-1}b_0}_1<\cdots<
       \overline{a_r^{-1}b_0}_1 .\] All displayed differences are in \(U\): the left terms are successive factors of the \(b\)-chain, the right terms are right factors of \(a_0^{-1}b_0\), and each mixed product equals \(a_i^{-1}b_j\). Applying the forward construction recovers the original simplex after left translation by \(a_0^{-1}\). The reverse composite similarly recovers the original \(E_T\) simplex. One-layer chains reduce to chains in \(U\). Deleting a vertex gives the same face, possibly with a different base vertex; that change is precisely a common left translation, which preserves all interval labels. Thus these simplex maps respect face identifications and descend to inverse cellular maps on the interval complexes.\par
This proves the unmarked identification used locally in Theorem 4.3. The formula is compatible with inclusions of spherical subsets. The global Salvetti homotopy conclusion still uses the compatible parabolic intersections and standard naturality-and-gluing result of [ps21]. Those are identified dependencies, not a newly established error in the source’s coincidence assertion. Asphericity of the Salvetti product implies that of the original Salvetti complex, as Corollary 4.4 requires.\par
\section*{Cyclic Artin groups: support and finite checks}
Lemmas 5.1–5.2 have the needed support content. A word equal to the identity cannot contain a unique occurrence of a simple generator. If \(s\notin\operatorname{Supp}(w)\) and \(sw=ws\), track its unique occurrence through braid moves between the reduced words \(sw,ws\). Before a noncommuting letter reaches its left, no first noncommuting braid involving \(s\) is possible. Thus \(s\) commutes with all of \(\operatorname{Supp}(w)\).\par
Lemma 5.3 needs a reducedness justification before strong exchange. The induction makes \(h=s_2\cdots s_n\cdots s_2\) reduced of length \(2n-3\); since \(s_1\) is outside its support, \(\ell(s_1h)=2n-2\). If \(s_1hs_1\) is nonreduced, its length is exactly \(2n-3\). The deletion word has this length and is reduced. Strong exchange and the singleton-letter argument then apply.\par
Lemma 5.4 first restricts to the connected union of the two reflection supports. Order reversal corresponds to inverting the product and reverses a possible counterexample. The generic intersection argument uses parabolic subgroups: a proper finite parabolic containing a specified codimension-one parabolic must equal it. This conclusion would not hold for arbitrary reflection subgroups. In the exceptional \(F_4\) case, the three full-support reflections of its \(B_3\) subgroup are \[s_1s_2s_3s_2s_1,\quad s_1s_3s_2s_3s_1,\quad
 s_2s_3s_2s_1s_2s_3s_2.\] Their displayed factorizations and conjugate-subgroup intersections were checked exactly; the decisive intersection orders are eight and two. Corollary 5.5 follows by selecting a reflection factor whose support crosses the endpoint of the predecessor-containing factor and applying ordered reflection rewriting [mcc15].\par
Exact arithmetic over \(\mathbb Q\) and \(\mathbb Q(\sqrt5)\) checked both path orientations for \(A_2\)–\(A_5\), \(B_2\)–\(B_5\), \(H_2,H_3,H_4,F_4\): 11,548 ordered reflection-pair implications, including 486 substantive predecessor/noncontainment candidates. All passed. These computations supplement the all-rank proof; finite sampling is not claimed to establish an infinite family.\par
Table 1 was checked analytically. For at least four vertices, labels at least six fail a proper three-vertex Gram test. Four vertices allow \(3333,3334,3335,3434,3435,3535\) up to symmetry; five allow \(33333,33334\). At least six vertices permit only three: a proper five-vertex path with an interior four already has nonpositive Gram determinant. Exact enumeration for four through eight vertices gives \(6,2,1,1,1\) classes. For the full cycle, the all-ones quadratic form is \(n-2\sum_e\cos(\pi/m_e)\le0\), proving nonsphericity. Connectedness gives irreducibility, including the affine cases.\par
In Proposition 5.6, omitted-vertex intervals intersect in the principal interval of their common spherical generators; their joins therefore agree on overlaps. In (4), membership of \(aw\) follows with length additivity. If \(a'=(au)\vee(av)=aww_0\), then \[|a|+|w|+|w_0|=|a'|
 \le |aw|+|w_0|\le |a|+|w|+|w_0|.\] Hence \(aw\le a'\le\delta_i\). Support containment alone would not prove this membership.\par
The proof of (6) has a definite but repairable factor-order error. Using the predecessor of \(\min I_u\) on \(au\) reverses Lemma 5.4’s left-factor hypothesis. Assume instead \(I_a\cup I_b=I_u\cup I_v=S\), and choose the predecessor of \(\min I_a\). It lies in \(I_u\) or \(I_v\). The correctly ordered lemma gives \(I_a\subset I_u\) or \(I_a\subset I_v\), contradicting \(bu\in U\) or \(bv\in U\), respectively. Thus one pair has proper support union and the required join. The cyclic conclusion survives this repair.\par
\section*{Combination and geometric consequences}
For Lemma 6.1, every oriented spherical Dynkin forest admits a linear extension. Compatible extensions differ by commuting swaps, hence give one Coxeter element. Deleting generators and reflection rewriting give nested intervals. For Lemma 6.2, cut the relevant tree at the edge between the selected support differences. A wrong orientation makes both \(r_1r_2\) and \(r_2r_1\) rank-two upper bounds in the finite interval and forces equality. Their positive root normals have a strictly negative inner product across the cut, so the independent reflections cannot commute. This also handles reducible ambient spherical diagrams.\par
Global atomic joins in Theorem 6.3 are computed in the finite carrier of the support union; principal-interval intersections make them independent of the carrier. For multiplying a join, the total support is complete and has the specified cyclic-by-spherical decomposition. The exact mixed case (2) on p.46 is empty: if only \(u\) lies in the cyclic factor, its entire nonspherical vertex set would equal \(\operatorname{Supp}(a_1)\cup\operatorname{Supp}(u)\), contained in the spherical support of \(au\in U\). An initially considered example using a larger ambient factor decomposition actually lies in the all-spherical case; it was rejected before assessment and is not a finding.\par
For the final rectangle, a complete same-side support union gives a join by a finite carrier or the cyclic criterion. Otherwise choose nonadjacent \(s,t\) in the first pair and nonadjacent \(s_u,s_v\) in the second. The four complete cross unions place these vertices in the exact support differences needed by Lemma 6.2. All cross edges have label two or point toward \(s_u,s_v\). The four-cycle hypothesis forces one of the two deliberately absent diagonals, a contradiction. This checks the distinctness and difference-set quantifiers in the final abbreviated argument.\par
Both Figure 1 examples were actually inspected and independently parsed from the source vector data. They have, respectively, seven vertices/seventeen edges/two maximal cliques and fourteen vertices/forty-two edges/six maximal cliques. All maximal cliques have the stated cyclic types, their commuting complements are empty, and neither graph has an induced four-cycle. These finite checks support the examples, not the general theorem.\par
Theorem 2.15 and the Bestvina paragraph need a centrality qualification: \(\langle\Delta\rangle\) need not be normal. In \(B_3\), classical \(\Delta=aba\) maps to a transposition changed by conjugation by \(a\). Normality, combined with exponent sum, would force \(a\Delta a^{-1}=\Delta\), a contradiction. Use a central power or assume centrality. The paper’s \(\Delta_E=(1,1)\) is already central.\par
For Theorem D’s generic direct-factor step, let \(H=G\times\mathbb Z\) be Garside, let \(d\) be a central Garside power and \(z\) generate the displayed split factor. For large \(N\), \(D=d^Nz\) is positive and central and contains the old simple divisors. Atomicity makes its divisor set finite; it is a Garside element for the same monoid. Choose \(N\) so its \(\mathbb Z\)-coordinate is nonzero. Then \(G\) injects with finite index into \(H/\langle D\rangle\), and the quotient geometric action restricts geometrically. Translation-number claims need the specified combinatorial metric, not arbitrary rescalings of arbitrary actions.\par
The external metric statements checked are injectivity for the Garside complex’s piecewise \(\ell^\infty\) orthoscheme metric [hae24] and the central-power CUB quotient statement [hae22]. Helly is asserted for \(A_\Gamma\times\mathbb Z\), CUB for \(A_\Gamma\); no general CAT(0) conclusion is proved. For the cyclic center corollary, the verified irreducibility and nonsphericity supply Jankiewicz–Schreve’s no-spherical-factor hypothesis after the checked \(K(\pi,1)\) argument [js23]. It does not extend to all joins with spherical factors.\par
\section*{Supporting checks}
The complete 2212-line extracted text, source-view receipts, independent counterexample scripts, exact Coxeter calculations, coverage and findings are retained with current artifact hashes. Finite calculations used independently written code and the approved Python runtime; no author code was obtained or executed.\par
External reading was bounded to the stated definitions, theorems and arguments: McCammond’s interval languages, Digne’s germ criterion, the published Soergel conditions, parabolic factor containment, reflection rewriting, and Paolini–Salvetti’s gluing argument. The complete global Garside normal-form theory and every general Farrell–Jones, Baum–Connes or semihyperbolicity dependency were not independently reproved. The Haettel injective-metric check used arXiv:2109.07891v5. The CUB central-power check used the later arXiv:2211.07857v2, Corollary 9.11, not an unread 2022 version. These dependency checks do not change the supplied version under review. Historical priority and the entire bibliography were not independently audited.\par
The lead worked with native in-chat helpers in disjoint directories. Their session self-description is GPT-6 Astra Pro; backend identity is not independently instrumented, and none is a human reviewer. The current TeX is compiled with the approved cached Tectonic runtime, untrusted mode and no shell escape. Exact command, compiler/dependency/log/TeX/PDF hashes, return code and actual report-page inspections are recorded against the current files. Missing optional software is not a mathematical defect. Scientific decision is pending; human reviews are empty and author adoption is unknown.\par
10 T. Haettel and J. Huang, New Garside structures and applications to Artin groups, arXiv:2305.11622v2 (2024). https://arxiv.org/abs/2305.11622v2. J. McCammond, An introduction to Garside structures, author manuscript, 15 February 2005, Definitions 1.3–1.16 and Theorems 1.7, 1.17. https://web.math.ucsb.edu/\textasciitilde{}jon.mccammond/papers/intro-garside.pdf. F. Digne, A Garside presentation for Artin–Tits groups of type affine C, Ann. Inst. Fourier 62 (2012), 641–666, Section 2. https://doi.org/10.5802/aif.2690. M. Soergel, Systolic complexes and group presentations, Groups Geom. Dyn. 17 (2023), 899–921, restricted-triangular definitions and Theorem 3.5. https://doi.org/10.4171/GGD/717. B. Baumeister, M. Dyer, C. Stump and P. Wegener, A note on the transitive Hurwitz action on decompositions of parabolic Coxeter elements, arXiv:1402.2500v1, Theorem 1.4. https://arxiv.org/abs/1402.2500v1. J. McCammond, Dual Euclidean Artin groups and the failure of the lattice property, author manuscript, 12 November 2014, Lemma 3.7. https://web.math.ucsb.edu/\textasciitilde{}jon.mccammond/papers/lattice-failure.pdf. G. Paolini and M. Salvetti, Proof of the K(pi,1) conjecture for affine Artin groups, Invent. Math. 224 (2021), 487–572, Section 5.2 and Theorem 5.5. https://doi.org/10.1007/s00222-020-01016-y. T. Haettel, Lattices, injective metrics and the K(pi,1) conjecture, arXiv:2109.07891v5, Theorem E and Corollary 3.12. https://arxiv.org/html/2109.07891v5. T. Haettel, A link condition for simplicial complexes, and CUB spaces, arXiv:2211.07857v2 (2025), Section 9.3, Corollary 9.11. https://arxiv.org/abs/2211.07857v2. K. Jankiewicz and K. Schreve, The K(pi,1)-conjecture implies the center conjecture for Artin groups, arXiv:2201.06591, Theorems 3 and 12. https://arxiv.org/abs/2201.06591.\par
\end{document}
